\documentclass[12pt]{article}
\usepackage[T1]{fontenc}
\usepackage[utf8]{inputenc}
\usepackage{lmodern}
\usepackage{textcomp}
\usepackage{enumitem}
\usepackage{geometry}
\geometry{margin=1in}
\begin{document}
\begin{center}
Answer Key - Chemistry - Undergraduate Level Assessment 19 \
\small (Answers are ordered exactly as in the question paper.)
\end{center}

\section*{Multiple Choice}
\begin{enumerate}[label=\arabic*.]
\item \textbf{Answer:} C
\\ \textit{Options:} A) K+; B) Ca$_{2}$+; C) Sc$_{3}$+; D) Rb+
\\ \textit{Note:} Sc$_{3}$+ has the smallest radius because it has lost three electrons, resulting in a higher effective nuclear charge that pulls the remaining electrons closer to the nucleus, leading to a smaller ionic radius compared to the other ions listed.
\item \textbf{Answer:} D
\\ \textit{Options:} A) Metallic element; B) Sulfide; C) Fluoride; D) Oxide
\\ \textit{Note:} Silicon is most commonly found in nature as silicon dioxide (SiO 2), which forms the basis of many minerals and rocks, making oxide the most prevalent naturally-occurring form.
\item \textbf{Answer:} D
\\ \textit{Options:} A) The orbitals are degenerate.; B) The set of orbitals has a tetrahedral geometry.; C) These orbitals are constructed from a linear combination of atomic orbitals.; D) Each hybrid orbital may hold four electrons.
\\ \textit{Note:} The statement is correct because each $sp^3$ hybrid orbital can hold a maximum of two electrons, not four, as per the Pauli exclusion principle.
\item \textbf{Answer:} A
\\ \textit{Options:} A) Cu; B) Gd; C) Eu; D) Am
\\ \textit{Note:} Copper (Cu) has a filled 3 d orbital and a partially filled 4 s orbital, but it does not have partially filled 4 f or 5 f orbitals, distinguishing it from elements that do.
\item \textbf{Answer:} A
\\ \textit{Options:} A) Blackbody radiation curves; B) Emission spectra of diatomic molecules; C) Electron diffraction patterns; D) Temperature dependence of reaction rates
\\ \textit{Note:} Planck's quantum theory successfully explained the observed blackbody radiation curves by introducing the concept of quantized energy levels, which resolved the ultraviolet catastrophe predicted by classical physics.
\end{enumerate}

\section*{True / False}
\begin{enumerate}[label=\arabic*.]
\setcounter{enumi}{5}
\item \textbf{Answer:} False
\\ \textit{Options:} A) True; B) False
\\ \textit{Note:} The statement is false because the magnetic field strength of 0.5 T is insufficient to induce significant polarization in 13 C at 298 K, as the polarization is highly dependent on the magnetic susceptibility and the temperature of the material.
\item \textbf{Answer:} True
\\ \textit{Options:} A) True; B) False
\\ \textit{Note:} The molecular geometry of thionyl chloride (SOCl 2) is trigonal pyramidal because it has a central sulfur atom bonded to two chlorine atoms and one oxygen atom, along with a lone pair of electrons, resulting in a geometry that is based on the tetrahedral arrangement of electron pairs.
\item \textbf{Answer:} True
\\ \textit{Options:} A) True; B) False
\\ \textit{Note:} Most organic functional groups have distinct vibrational modes that correspond to specific wavelengths in the infrared spectrum, allowing for their identification and aiding in structural analysis through infrared spectroscopy.
\item \textbf{Answer:} True
\\ \textit{Options:} A) True; B) False
\\ \textit{Note:} The amide ion, NH 2-, is considered the strongest base in liquid ammonia because it can readily accept protons, resulting in the formation of ammonia and demonstrating a high affinity for protons in this solvent system.
\item \textbf{Answer:} False
\\ \textit{Options:} A) True; B) False
\\ \textit{Note:} The statement is false because hexane has only two distinct chemical shifts in its 13 C NMR spectrum due to its symmetrical structure, which results in equivalent carbon environments.
\end{enumerate}

\section*{Long Form Response}
\begin{enumerate}[label=\arabic*.]
\setcounter{enumi}{10}
\item \textbf{Answer:} Cobalt-60 is produced from cobalt-59 through a process known as neutron bombardment, where cobalt-59 nuclei are exposed to neutrons in a nuclear reactor or particle accelerator. During this interaction, cobalt-59 captures a neutron and undergoes beta decay, transforming into cobalt-60, which is a radioactive isotope. The significance of cobalt-60 in radiation therapy lies in its ability to emit high-energy gamma rays when it decays, which can effectively target and destroy cancer cells while minimizing damage to surrounding healthy tissue. This property makes cobalt-60 an invaluable source of radiation for external beam radiotherapy, allowing for precise treatment that enhances patient outcomes in cancer management.
\\ \textit{Note:} correct because it accurately describes the neutron bombardment process that converts stable cobalt-59 into radioactive cobalt-60, and it highlights the therapeutic significance of cobalt-60's gamma radiation in effectively targeting and destroying cancer cells in radiation therapy.
\item \textbf{Answer:} The Electron Paramagnetic Resonance (EPR) spectrum of a rigid nitronyl nitroxide diradical, characterized by a coupling constant J that is much smaller than the hyperfine coupling constant a, reveals specific features due to the interaction of the unpaired electrons with the nitrogen nuclei. In this scenario, the nitrogen nuclei, which possess a nuclear spin, couple with the unpaired electrons, leading to a splitting of the energy levels. As a result, this coupling contributes to the formation of multiple spectral lines in the EPR spectrum. The number of lines produced is determined by the multiplicity of the nitrogen nuclei, which increases the complexity of the spectrum. Specifically, for a system with two equivalent nitrogen nuclei, one can expect to observe a pattern of six lines, reflecting the various spin states resulting from this coupling. Overall, this interaction enhances our understanding of the electronic environment and dynamics within the diradical.
\\ \textit{Note:} The answer correctly highlights that in a rigid nitronyl nitroxide diradical with J \textless{}\textless{} a, the interaction between the unpaired electrons and nitrogen nuclei leads to energy level splitting, which directly influences the number of spectral lines observed in the EPR spectrum.
\end{enumerate}

\vfill
\noindent\textit{End of Answer Key}
\end{document}